\section{Limitations}
\label{sec:limitations}

Fine-teacher predictions are model-derived soft supervision rather than human
ground truth, and both gate corpora are likewise model-derived.
Our guards therefore constrain a checkpoint without certifying it.
The one held-out set with human labels is flat across every variant we
measured, and released M7 remains ahead there, so the frozen benchmark
establishes an ordering under its own supervision rather than a general one.

Milestone-to-milestone variation is comparable to the differences that
selection rules are asked to resolve.
We measured seed noise with one full-horizon replicate and report anything
inside that band as unresolved, but a single replicate is a weak estimate of a
variance, and we have not repeated the shipped configuration a third time.

The ablations in Appendix~\ref{sec:also-tried} are negative results measured
under one corpus, one optimizer, and one horizon.
They establish that those terms do not earn their place in this configuration.
They do not establish that no version of a medial, separation, or connectivity
objective could help under a different corpus or schedule, and the medial tail
floor in particular remains untested at scale because the corpus it needs does
not exist in the required form.

Test-time augmentation improves the shipped number at eight times the inference
cost.
We report both the augmented and unaugmented figures because that trade is a
deployment decision, not a property of the model.

The teaser shows the shipped model; the remaining rendered figures still show
the earlier v31 student and have not been regenerated against the shipped
weights (Sec.~\ref{sec:figure-provenance}).
The blob audit of Sec.~\ref{sec:blob-audit} has no ground truth: solid-mass
share is a geometric proxy, and it cannot see two wraps fused where the CT
itself shows no gap. The HercUNet fidelity check covers four patches, not the
whole benchmark.
The student is not free of masses either: in some open cubes it fuses wraps
itself.
On the seeded-random open PHerc0191 cube shown in the release record,
$\openFuseFz\%$ of its predicted voxels lie inside a solid mass, against
$\openFuseMSeven\%$ for the published M7.
The downstream postprocessors of Appendix~\ref{sec:fitter} keep provenance pins
at the earlier operating threshold and must be re-qualified before they are
combined with the shipped inference recipe.

The curve fitter has a different failure mode from the model.
A completion may be locally smooth and persistent across slices but still join
the wrong wrap near the core.
Radial, third-component, and cross-sheet gates reduce this risk but cannot prove
global correctness.
The postprocess must remain independently measured and its additions must remain
recoverable during qualification.

\begin{figure*}[t]
  \centering
  \figureasset{figure_4_failure_cases}{0.96\linewidth}{2.6in}
  \caption{Known edge cases: a component-count proxy that disagrees with the
  visible structure, mild drift or thinning on locked sheets, and blind
  PHerc1447 views that remain undergrown or fragmented.
  Every row includes CT, released M7, the available model-derived reference,
  the raw student, and a two-sided comparison. Yellow reference masks are not
  ground truth. Student panels show the v31 configuration; see
  Sec.~\ref{sec:figure-provenance}.}
  \Description{Edge cases showing topology-metric disagreement, mild drift or
  thinning, and residual undergrowth.}
  \label{fig:limitations}
\end{figure*}

Finally, the shipped checkpoint and the code are public
(Sec.~\ref{sec:release}), but exact replay of training and evaluation still
requires large, separately licensed artifacts.
The validation corpus, teacher atlases, M7 weights, and the unshipped milestone
checkpoints require dedicated artifact hosting and resolved upstream terms.
